% Lösungen Einheit 1 von 2 – Genau k Treffer ohne Zurücklegen: die Situation erkennen (feste kleine Gesamtheit, Ziehen ohne Zurücklegen (zusammenbau v0.1)
\einheitenkopf{Einheit 1 von 2 · Genau k Treffer ohne Zurücklegen: die Situation erkennen (feste kleine Gesamtheit, Ziehen ohne Zurücklegen}

\erg{7}{a) ohne Zurücklegen – Auswahlen zählen: niemand wird zweimal gelost, die Klasse schrumpft. \quad b) $1 - P(\text{alle weiß}) = 1 - \frac{5}{8} \cdot \frac{4}{7} \cdot \frac{3}{6} = 1 - \frac{5}{28} = \frac{23}{28}$ \quad c) $P = \frac{(9 \text{ über } 2) \cdot (51 \text{ über } 6)}{(60 \text{ über } 8)} \approx 0{,}2534$, also etwa $25{,}3\,\%$. Binomial mit $p = 0{,}15$ käme etwa $0{,}2376$ heraus – falsch, weil ohne Zurücklegen ausgewählt wird. \quad d) $P(X \le 3) = 1 - P(X = 4) - P(X = 5) = 1 - \frac{70 \cdot 12}{15\,504} - \frac{56}{15\,504} \approx 0{,}9422$ \quad e) $P = 1 - \frac{(119 \text{ über } 4)}{(120 \text{ über } 4)} = 1 - \frac{116}{120} = \frac{1}{30} \approx 0{,}0333$ – jede der $120$ Personen hat dieselbe Chance, unter den $4$ Gezogenen zu sein. Binomial ungeeignet: es wird ohne Zurücklegen gelost, die Ziehungen sind nicht unabhängig und die Trefferwahrscheinlichkeit bleibt nicht gleich.}
\erg{8}{a) Fehler: Mia rechnet binomial mit fester Wahrscheinlichkeit $0{,}2$, obwohl ohne Zurücklegen gezogen wird. Richtig: $P = \frac{(4 \text{ über } 1) \cdot (16 \text{ über } 2)}{(20 \text{ über } 3)} = \frac{480}{1140} = \frac{8}{19} \approx 0{,}4211$ \quad b) Jede gezogene Kugel fehlt danach in der Urne: Gesamtzahl und – bei einem Treffer – auch die Zahl der Treffer werden kleiner. Damit ändert sich der Anteil der Treffer und so die Wahrscheinlichkeit beim nächsten Zug. \quad c) $P(\text{Ablehnung}) = 1 - \frac{(47 \text{ über } 5)}{(50 \text{ über } 5)} \approx 1 - 0{,}7240 = 0{,}2760$ – gut jede vierte solche Lieferung wird abgelehnt.}
